\section[From spectral filters to geometric messages]
{A second longitudinal case: from spectral filters to geometric messages}

The phrase ``graph convolutional network'' now covers a lineage of spectral filtering,
neighborhood aggregation, and typed message passing.  Spectral networks begin from functional
calculus of a graph Laplacian.  The
Kipf--Welling GCN retains a first-order normalized propagation rule.  GraphSAGE turns propagation
into permutation-invariant neighborhood aggregation.  GAT makes the route state-dependent.  MPNNs
allow typed edge messages, GIN asks which neighborhood multisets survive aggregation, and geometric
networks constrain how scalar, vector, and coordinate states transform.  Each transition changes
what an edge does as well as increasing capacity.

This history is especially useful here because the equations look continuous while their meanings
change.  The audit will follow
\begin{equation}
\boxed{
\begin{aligned}
\text{spectral function}&\longrightarrow\text{local polynomial},\\
\text{fixed normalized route}&\longrightarrow\text{adaptive aggregation},\\
&\longrightarrow\text{typed geometric message}.
\end{aligned}}
\label{eq:book-gcn-lineage}
\end{equation}
At no point does successful propagation alone identify a statistical interaction, a physical
force, or a biological mechanism.

\subsection{Spectral convolution is functional calculus on a declared graph}

Let $A$ be a symmetric nonnegative adjacency matrix, $D_{ii}=\sum_jA_{ij}$, and
\begin{equation}
L=I-D^{-1/2}AD^{-1/2}=U\Lambda U^{\mathsf T}
\label{eq:book-gcn-laplacian}
\end{equation}
be the normalized graph Laplacian.  The columns of $U$ replace Fourier modes: $U^{\mathsf T}x$
are the graph-frequency coefficients of a node signal $x$.  A spectral convolution is defined by
\begin{equation}
g_\theta *_G x
=Ug_\theta(\Lambda)U^{\mathsf T}x
=g_\theta(L)x.
\label{eq:book-gcn-spectral-convolution}
\end{equation}
The equality on the right is the important one.  The native object is a spectral function of a
specified Laplacian.  Low values of $\lambda$ describe signals that vary slowly across heavily
weighted edges, so a filter concentrated at low graph frequencies acts as a graph-relative
smoother \cite{bruna2014,defferrard2016}.

Nothing in Eq.~\ref{eq:book-gcn-spectral-convolution} discovers what the graph means.  If $A$ is a
covalent graph, a spatial contact graph, a sequence-neighborhood graph, or a learned similarity
graph, the same functional calculus has different semantics.  The spectrum belongs to the pair
$(L,\mathcal H_{\mathrm{node}})$ after the node state and edge weights have been declared.

The original spectral parameterization also has two practical liabilities.  Applying an arbitrary
$g_\theta(\Lambda)$ appears to require eigendecomposition, and parameters attached to the
eigenvalues of one graph transfer only through an additional alignment or parameterization.
Locality appears when the multiplier is restricted, since an unrestricted spectral multiplier
generally produces a dense operator in node coordinates.

\subsection{Chebyshev filters convert frequency response into finite-hop propagation}

ChebNet replaces the unrestricted spectral multiplier by a polynomial.  Rescale the spectrum to
$[-1,1]$ with
\begin{equation}
\widetilde L=\frac{2}{\lambda_{\max}}L-I
\end{equation}
and use Chebyshev polynomials $T_k$:
\begin{equation}
g_\theta(L)x
\approx
\sum_{k=0}^{K}\theta_kT_k(\widetilde L)x.
\label{eq:book-gcn-chebyshev}
\end{equation}
Because a degree-$k$ polynomial in $L$ connects only nodes reachable within $k$ graph steps, a
degree-$K$ filter is $K$-hop localized.  Spectral approximation has become a spatial propagation
rule without ceasing to be a function of the Laplacian.

This step exposes an assumption that the Fourier notation can hide.  A polynomial filter shares
coefficients across edges: the graph supplies the support and geometry, while the
shared coefficients $\theta_0,\ldots,\theta_K$ specify how different path lengths and spectral
scales are combined.  Two edges with the same local role are processed by the same polynomial.

\subsection{The classical GCN is a tied first-order filter}

The familiar GCN layer follows from a strong truncation of
Eq.~\ref{eq:book-gcn-chebyshev} \cite{kipf2017}.  Take $K=1$ and approximate
$\lambda_{\max}\simeq2$.  Since
\begin{equation}
\widetilde L=L-I=-D^{-1/2}AD^{-1/2},
\end{equation}
the one-channel filter becomes
\begin{equation}
\theta_0x-\theta_1D^{-1/2}AD^{-1/2}x.
\end{equation}
Tying $\theta_0=-\theta_1=\theta$ gives
\begin{equation}
\theta\left(I+D^{-1/2}AD^{-1/2}\right)x.
\label{eq:book-gcn-first-order}
\end{equation}
The subsequent renormalization trick adds self-loops before normalization:
\begin{equation}
\widehat A=A+I,
\qquad
\widehat D_{ii}=\sum_j\widehat A_{ij},
\qquad
S=\widehat D^{-1/2}\widehat A\widehat D^{-1/2}.
\label{eq:book-gcn-renormalized-route}
\end{equation}
For a feature matrix $H^{(\ell)}$, the layer is
\begin{equation}
H^{(\ell+1)}
=\sigma\!\left(SH^{(\ell)}W^{(\ell)}\right).
\label{eq:book-gcn-layer}
\end{equation}

Equation~\ref{eq:book-gcn-layer} should be read as two operators with different axes:
\begin{equation}
H
\xrightarrow{\;S\;}
\text{mix node states along the declared graph}
\xrightarrow{\;W\;}
\text{mix feature channels}.
\label{eq:book-gcn-two-operators}
\end{equation}
The propagation matrix $S$ is fixed once the input graph is fixed; the trainable matrix $W$ is
shared over nodes.  Standard GCN therefore learns channel transformations around a prescribed,
degree-normalized node-mixing rule, with edgewise freedom limited by the input graph.

The self-loop has a precise computational purpose: a node's current state participates in the next
state.  Its meaning is computational.  Degree normalization controls the relative scale of messages
from differently connected nodes.  This normalization belongs to the propagation geometry;
probability normalization over molecular conformations is a separate object, and the rows of the
symmetric normalization generally have varying sums.

\subsection{Depth reveals the diffusion assumption}

Ignoring nonlinearities and channel maps for a moment, stacking layers repeatedly applies $S$:
\begin{equation}
H^{(K)}\approx S^KH^{(0)}.
\label{eq:book-gcn-repeated-propagation}
\end{equation}
If $S=U\operatorname{diag}(\mu_r)U^{\mathsf T}$, then
\begin{equation}
S^K=U\operatorname{diag}(\mu_r^K)U^{\mathsf T}.
\label{eq:book-gcn-power-spectrum}
\end{equation}
Modes with $|\mu_r|<1$ decay, while dominant smooth modes survive.  Under the usual connectivity
and aperiodicity conditions, node representations within a component become increasingly aligned
with a small stationary subspace.  This is the linear skeleton of over-smoothing.  Nonlinearities,
learned channels, and residual paths alter the exact limit, while the spectral pressure created by
repeated propagation remains.

The phenomenon reveals the inductive bias encoded by depth.  GCN assumes that
neighbors should exchange and often homogenize information.  That is useful under homophily.  If
an edge instead connects complementary types, opposite labels, or directed roles, indiscriminate
low-pass smoothing can destroy the signal of interest.  Adjacency grants permission to
communicate; similarity requires an additional assumption about the signal and edge semantics.

APPNP makes the retained source explicit through a restart term \cite{klicpera2019}:
\begin{equation}
H^{(k+1)}=(1-\alpha)SH^{(k)}+\alpha H^{(0)}.
\label{eq:book-gcn-appnp-recursion}
\end{equation}
At convergence, when the inverse exists,
\begin{equation}
H^*=\alpha\left[I-(1-\alpha)S\right]^{-1}H^{(0)}.
\label{eq:book-gcn-appnp-resolvent}
\end{equation}
The result is a resolvent or personalized-PageRank filter with a persistent contribution from the
initial signal.  GCNII similarly preserves initial features and identity mappings across depth
\cite{chen2020gcnii}.  Both methods reshape how multi-hop evidence accumulates.

\subsection{GraphSAGE turns propagation into a set function}

GraphSAGE writes the update without committing to one full-graph matrix multiplication
\cite{hamilton2017}:
\begin{align}
m_i^{(\ell)}
&=\operatorname{AGG}^{(\ell)}
\left(\{h_j^{(\ell)}:j\in\mathcal N(i)\}\right),\nonumber\\
h_i^{(\ell+1)}
&=\sigma\!\left(
W^{(\ell)}[h_i^{(\ell)}\Vert m_i^{(\ell)}]
\right).
\label{eq:book-gcn-graphsage}
\end{align}
The aggregation must be invariant to an arbitrary enumeration of the neighbors.  Mean, sum,
pooling, and recurrent aggregators retain different information about the neighborhood multiset.
Neighbor sampling makes the computation inductive and scalable, but also turns exact propagation
into a sampling estimator.

The mathematical question has shifted.  Spectral GCN asks which function of a graph operator acts
on the signal.  GraphSAGE asks which permutation-invariant function summarizes a local multiset.
Both propagate over edges, but their approximation errors and information bottlenecks are typed
differently.

\subsection{GAT makes the route depend on the current state}

Graph attention replaces degree-determined coefficients by content-dependent scores
\cite{velickovic2018}:
\begin{align}
e_{ij}^{(\ell)}
&=a_\theta\!\left(W h_i^{(\ell)},Wh_j^{(\ell)},e_{ij}\right),\nonumber\\
\alpha_{ij}^{(\ell)}
&=\frac{\exp e_{ij}^{(\ell)}}
{\sum_{k\in\mathcal N(i)}\exp e_{ik}^{(\ell)}},\nonumber\\
h_i^{(\ell+1)}
&=\sigma\!\left(\sum_{j\in\mathcal N(i)}
\alpha_{ij}^{(\ell)}Wh_j^{(\ell)}\right).
\label{eq:book-gcn-gat}
\end{align}
Now the graph support may remain fixed while the weighted route changes with input and layer:
\begin{equation}
S\quad\longrightarrow\quad P_\theta(H,E).
\end{equation}
For each receiver $i$, the coefficients form a categorical distribution over allowed sources.  The
normalization makes $P_\theta$ row-stochastic under the receiver-to-source convention and generally
asymmetric.  These coefficients type the computational route; conditional dependence, interaction
energy, and physical force require separate extraction and validation.  The same distinction reappears in
Transformer attention on a dense token graph.

GAT replaces the fixed self-adjoint graph-relative filter of a GCN with an input-dependent,
generally directed propagation rule.
The appropriate spectral geometry, if one is needed, must be reconstructed for that directed
state-dependent operator.

\subsection{MPNNs promote the edge from a coefficient to a typed computation}

The message-passing neural network abstraction separates message, aggregation, and update
\cite{gilmer2017,battaglia2018relational}:
\begin{align}
m_{ij}^{(\ell)}
&=\phi_m^{(\ell)}
\left(h_i^{(\ell)},h_j^{(\ell)},e_{ij}\right),\nonumber\\
m_i^{(\ell)}
&=\bigoplus_{j\in\mathcal N(i)}m_{ij}^{(\ell)},\nonumber\\
h_i^{(\ell+1)}
&=\phi_u^{(\ell)}\left(h_i^{(\ell)},m_i^{(\ell)}\right).
\label{eq:book-gcn-mpnn}
\end{align}
Classical GCN is recovered by choosing a linear isotropic message such as
$m_{ij}=S_{ij}h_jW$.  General MPNNs can condition on bond type, distance, orientation, molecular
identity, or a hidden edge vector.  The edge is no longer exhausted by one scalar coefficient; it
may implement a map between typed node fibers.

Depth also changes interaction order.  A two-layer update at $i$ can depend on paths $k\to j\to i$
and therefore on configurations of several nodes even when every primitive message is pair-indexed.
Pairwise execution support can compose into higher-order output dependence across layers.

GIN isolates a separate issue: whether aggregation distinguishes different neighborhood
multisets \cite{xu2019gin}.  Its update is
\begin{equation}
h_i^{(\ell+1)}
=\operatorname{MLP}^{(\ell)}\!\left(
(1+\epsilon^{(\ell)})h_i^{(\ell)}
+\sum_{j\in\mathcal N(i)}h_j^{(\ell)}
\right).
\label{eq:book-gcn-gin}
\end{equation}
With appropriate injective maps, it reaches the discriminative power of the one-dimensional
Weisfeiler--Lehman test.  This is an expressivity statement relative to a combinatorial test, not a
claim that the recovered features are causal or physical.  Mean aggregation, sum aggregation, and
scalar edge compression can agree on a downstream benchmark while discarding different structural
information.

\subsection{Geometric networks add transformation law, not molecular energy}

For molecular systems the node state may include coordinates $r_i\in\mathbb R^3$.  A simple
invariant message can depend on squared distance,
\begin{equation}
m_{ij}=\phi_m(h_i,h_j,e_{ij},\|r_i-r_j\|^2),
\end{equation}
while an equivariant coordinate update has the schematic form
\begin{equation}
r_i'=r_i+\sum_j(r_i-r_j)\phi_x(m_{ij}).
\label{eq:book-gcn-equivariant-update}
\end{equation}
Under a rigid transformation $r_i\mapsto Qr_i+t$, invariant scalar messages remain unchanged and
the displacement in Eq.~\ref{eq:book-gcn-equivariant-update} transforms by $Q$.  Consequently
\begin{equation}
F(QR+t)=QF(R)+t
\label{eq:book-gcn-equivariance-law}
\end{equation}
for the coordinate map under the stated architecture \cite{satorras2021}.  More general geometric
networks carry vectors, tensors, or irreducible rotation representations and type every message by
its transformation law.

Equivariance licenses a symmetry statement.  Energy, momentum conservation, and physical time
evolution additionally require integrability, conservation laws, units, and a dynamical model.
The same distinction applied to invariant point attention and conditional diffusion in
Chapter~8.

\subsection{Two diffusions that share a word but not a state space}

Repeated GCN propagation and a coordinate diffusion model are easily conflated because both use
diffusion language.  Their operators act on different objects:
\begin{equation}
\begin{aligned}
\text{graph-signal diffusion:}&\quad
H^{(k+1)}=SH^{(k)},
\quad S:\mathcal H_{\mathrm{node}}\to\mathcal H_{\mathrm{node}},\\
\text{generative diffusion:}&\quad
\mathrm dR_t=f(R_t,t)\,\mathrm dt+g(t)\,\mathrm dW_t,\\
&\quad p_t\text{ is a density on coordinate configurations}.
\end{aligned}
\label{eq:book-gcn-two-diffusions}
\end{equation}
The first smooths features relative to a declared graph.  The second smooths a probability density
in a configuration space and learns its score for reverse transport.  A graph Laplacian can appear
inside a generative model.  The state space and generator determine which diffusion is meant and
which physical interpretation is available.

\subsection{What actually changed across the lineage}

The longitudinal compression is
\begin{equation}
\boxed{
\begin{aligned}
g_\theta(L)&\longrightarrow\sum_k\theta_kL^k\longrightarrow SHW,\\
&\longrightarrow P_\theta(H)HW
\longrightarrow\phi_m(h_i,h_j,e_{ij}),\\
&\longrightarrow\text{typed equivariant edge maps}.
\end{aligned}}
\label{eq:book-gcn-operator-progression}
\end{equation}
The progression relocates complexity.  Spectral GCN places it in the graph filter; classical GCN
compresses it into a fixed normalized route plus learned channel maps; attention moves some of it
into state-dependent routing; MPNNs place it inside message functions; geometric networks add
exact transformation constraints.  None of these moves monotonically increases interpretability.

For protein models, this history supplies a clean audit.  Is the adjacency observed, constructed,
or generated from the current state?  Is an edge a scalar route, a hidden vector, or an operator
between typed fibers?  Does depth merely spread evidence, or does it create a higher-order output
response?  Is symmetry guaranteed by architecture, and what physical claim remains unlicensed?
AlphaFold's persistent pair track and triangle computations go beyond a classical node-only GCN,
but they inherit the same central design question: which relations deserve state, and what may be
concluded from the way that state is propagated?

\subsection{One audit ledger across the lineage}

The lineage can now be returned to the same coordinates used for the model portraits:
\begin{center}
\small
\begin{tabular}{@{}>{\raggedright\arraybackslash}p{0.17\linewidth}
                >{\raggedright\arraybackslash}p{0.24\linewidth}
                >{\raggedright\arraybackslash}p{0.25\linewidth}
                >{\raggedright\arraybackslash}p{0.24\linewidth}@{}}
\toprule
Family & State, support, and edge carrier & Context, dynamics, and normalization & Estimator and first unlicensed claim \\
\midrule
Spectral/ChebNet & node signals on a declared graph; edge meaning inherited from $A$ & fixed self-adjoint filter; polynomial locality & supervised filter coefficients; physical meaning remains inherited from the declared graph \\
GCN/APPNP & node features; fixed degree-normalized scalar route & repeated smoothing, with optional restart to initial state & task loss; propagation defines computational smoothing, with statistical or molecular meaning requiring a corresponding generative law \\
GraphSAGE and GAT & sampled neighborhoods or state-dependent routes & set aggregation or row-normalized adaptive routing & task loss and neighbor sampling; attention weights describe adaptive routes \\
MPNN, GIN, and geometric networks & typed node and edge states, optionally coordinates & compositional messages; equivariant variants constrain transformation laws & task-specific loss; force, energy, and mechanism claims additionally require their defining physical structure \\
\bottomrule
\end{tabular}
\end{center}

The common entry is computation on a graph.  What changes is the graph's provenance, the object
carried by an edge, the information available to the message, and the transformation law of the
state.  Those are precisely the coordinates in Eq.~\ref{eq:model-audit-map}; together they determine
the meaning carried by each historical name.
